\chapter{Dialogs and Form Workflows}
\label{ch:dialogs-forms}
\chapterquote{A dialog is a temporary window around a durable decision. Keep those responsibilities separate.}{The dialog rule}

\begin{chaptermap}
Core question & How should form dialogs coordinate commit, cancel, validation, keyboard behavior, dirty state, and cleanup without becoming piles of listeners?\\
Primary APIs & \api{FormDialog}, \api{FormDialogShell}, \api{FormDialogLifecycle}, \api{FormDialogValidationBinding}, \api{FormDialogKeyboardBinding}, \api{ProblemFocusResolver}, \api{SwingEditors}, \api{DirtyState}, \api{DirtyStates}, \api{FormDialogActionState}, \api{RevertPolicy}, and \api{CancelConfirmation}.\\
Layout topic & MigLayout dialog bodies with generated or handwritten field bindings.\\
Corusco connection & Corusco centralizes dialog semantics while leaving component layout and business workflow in application code.\\
Companion examples & \dialogExample, \dialogApplyRevertExample, \dialogValidationExample, \dialogActiveEditorExample, \dialogDirtyCancelExample, \dialogLifecycleExample, \dialogKeyboardExample, and \multiFormDialogExample.\\
\end{chaptermap}

\section{The modal transaction}

\termdef{Modal dialog}{in Swing} is a window that temporarily prevents the user
from interacting with some other window until the dialog reaches a terminal
decision. The word terminal is important. A modal dialog is not merely a
\api{JDialog} with fields in it. It is a small transaction in which the program
asks the user for a value, a decision, a correction, or a confirmation, and then
returns control to the surrounding workflow. The transaction may succeed,
cancel, revert, or fail to accept the current input. It should not half-succeed
by closing the window while leaving the caller to guess what happened.

\termdef{Form model}{in this book} is the object that owns editable state,
validation problems, reset, dirty-baseline acceptance, and result creation for a
form. It may be a handwritten \api{FormModel}, a generated model, or a facade
over richer domain state. The key property is that the durable edit state is not
stored only in the Swing widgets. Text fields, combo boxes, tables, and custom
editors are user-interface instruments; the form model is the object that can
answer ``what would be committed if OK succeeded now?'' and ``what is wrong
with the current input?''

\termdef{Dialog controller}{in Corusco Swing} is an object that coordinates the
standard modal rules without being a native window. \api{FormDialog} is such a
controller. It owns OK, Apply, Revert, and Cancel semantics for a form model. It does
not build the form body, it does not show a native \api{JDialog}, and it does
not decide which repository method will save the result. That restraint is the
point. A controller that owns only the decision protocol can be tested without a
screen and reused without imposing a visual form language.

\termdef{Dialog shell}{in this chapter} is the temporary Swing window that hosts
the form body and the button row. A shell may be a \api{JDialog}, an internal
frame, an application-specific modal service, or a test harness that never
shows a native peer. The shell owns modality, ownership, placement, focus
activation, and disposal. The shell should not be the only place where the
dialog's business result lives.

\begin{lstlisting}[caption={Controller first, native shell second.}]
FormDialog<CustomerEditForm, CustomerEdit> controller =
        new FormDialog<>(form, rootPanel, dirtyState, cancelConfirmation);

FormDialogShell<CustomerEditForm, CustomerEdit> shell =
        FormDialogShell.create(owner, "Edit Customer", controller);

DialogResult<CustomerEdit> result = shell.showModal();
\end{lstlisting}

The exact shell helper may vary by application, but the separation should not.
The form model owns editable facts. The root panel owns Swing components and
layout. Bindings connect the two. The controller owns the modal decision. The
shell owns the window. The caller owns what to do with an accepted value. When a
program merges all five roles into one class named \api{CustomerDialog}, it
usually starts pleasantly and ends with an untestable knot: private booleans for
dirty state, private fields for results, action listeners that update labels,
anonymous keyboard actions, and cleanup code that runs only on the branch that
the author happened to remember.

This is not a complaint about \api{JDialog}. Swing gives us a useful native
window. The problem is that \api{JDialog} is a container, not a workflow model.
It knows how to be shown, hidden, packed, placed, and disposed. It does not know
that a formatted text field may still contain an uncommitted parse candidate,
that Apply should move the dirty baseline but not close the window, that Escape
should respect a dirty-cancel policy, or that validation focus may need to
select a hidden tab before requesting focus. Those rules belong to the
application architecture around the shell.

\begin{figure}[htbp]
\centering
\begin{minipage}[t]{0.48\linewidth}
\includegraphics[width=\linewidth]{\dialogValidationFigure}
\end{minipage}\hfill
\begin{minipage}[t]{0.48\linewidth}
\includegraphics[width=\linewidth]{\dialogActiveEditorFigure}
\end{minipage}
\caption{Dialog workflow has visible states. Validation should keep the window open and focusable, while OK must commit active editors before reading the model.}
\label{fig:dialog-form-screens}
\end{figure}

The two states in Figure~\ref{fig:dialog-form-screens} look ordinary because
they are ordinary. A useful dialog spends most of its time in ordinary states:
editing, invalid, dirty, busy, applying, or waiting for a confirmation. The
trouble is that Swing makes each state easy to implement once and easy to
forget in the next dialog. A table editor has a last edit; a text field has a
parse state; a validation summary has a most severe problem; a cancel button
has a dirty policy; a root pane has a default button and an Escape mapping; a
background operation has a cancellation boundary. Corusco's dialog classes are
small because they do not try to make any of these concerns exotic. They give
the repeated concerns stable names and a single execution order.

There is a useful mental model from database transactions, but it must be used
carefully. A dialog is transactional in the user-interface sense: it presents a
temporary editing context and returns a result only after acceptance. It is not
usually a database transaction. Holding a database transaction open while the
user thinks, switches windows, takes a phone call, or edits a table is normally
wrong. The dialog should return a result to the application workflow; the
workflow may then start an application-level save, possibly asynchronously,
possibly with conflict handling. Closing the modal shell should not by itself
mean that the outside world has been changed.

For this reason the first design question for a dialog is not ``which buttons
will it have?'' but ``what is the result protocol?'' The protocol names the
accepted value, the cancelled state, the state that prevents acceptance, and
the cleanup that follows either terminal outcome. A one-field rename dialog
still has such a protocol, although it may be tiny. A multi-tab customer editor
has the same protocol, although it must coordinate more fields, validators,
tables, and background checks. Once the protocol is visible, the button row is
just one way of invoking it.

The second design question is ``where does incomplete input live?'' In a poor
dialog it lives partly in text fields, partly in table editors, partly in
temporary local variables, and partly in a domain object that has already been
mutated. In a disciplined dialog the incomplete input has one presentation
model. Widgets can contain transient text during editing, but the form model is
the place where the dialog evaluates committability and creates its result.
This keeps invalid user input from leaking into committed domain objects and
keeps the view from becoming the only source of truth.

The third design question is ``who owns temporary relationships?'' A dialog
activation installs many short-lived relationships: document listeners,
selection listeners, command adapters, validation decorations, keyboard
bindings, task callbacks, popup helpers, generated behavior plans, and perhaps
a cache attached to a presenter. These things should end with the dialog
activation. If the dialog object is reused, they should end between
activations. If the owner window is disposed, they should end even if the user
never pressed Cancel. Corusco's lifecycle objects exist because these temporary
relationships are where many Swing memory leaks are born.

It is useful to classify dialogs by the kind of decision they return. A data
entry dialog returns a newly entered value. An edit dialog returns a replacement
or patch for an existing value. A chooser returns one value from a larger set. A
confirmation returns permission to continue. A property sheet may return
several applied snapshots before it finally closes. A wizard is a sequence of
dialogs hidden behind one shell; each page may be invalid, dirty, or busy even
though the final result is produced only at the end. These categories look
different on the screen, but they all need an explicit result protocol. The
protocol says when the user has produced a value, when cancellation is allowed,
and which temporary resources must be released.

The classification prevents two common mistakes. The first mistake is treating
every dialog as a data-entry editor. A confirmation dialog should not have a
form model just because the data-entry editor has one; it may need only a
command and a result. The second mistake is treating every editor as a custom
case. A chooser with a filter field still has dirty or pending state. A wizard
page still has validation problems and focus recovery. A property sheet still
has active editors. If the shared protocol is visible, the special parts stay
small and the common rules remain common.

The same classification also clarifies ownership of domain mutations. A chooser
normally should not mutate the chosen object. A data-entry dialog normally
returns a value to be inserted elsewhere. An edit dialog may edit a detached
copy, a patch object, or a presentation model that later updates the domain.
A property sheet with Apply may mutate the live application state after each
successful Apply. These choices affect dirty baselines and undo policy. They
should be decided before the first button listener is written, because the
listener will otherwise smuggle in the policy one branch at a time.

\begin{coruscobox}{Explicit Modal Protocols}
Plain Swing gives excellent building blocks for modal windows, but it leaves
the modal protocol implicit. Corusco makes the protocol explicit: a
\api{FormDialog} coordinates commit, validation, result creation, Apply,
Cancel, dirty confirmation, reset, command enablement, and one-shot close; the
native shell remains ordinary Swing.
\end{coruscobox}

\section{OK, Apply, Cancel, and Close}

The OK button does not mean ``close now.'' It means ``attempt to produce a
valid result.'' The distinction is not philosophical; it is the difference
between preserving and losing user data. A correct OK path first asks active
editors to publish their values, then refreshes committability, then reads the
form result, then accepts the current values as the new baseline, then records
an accepted dialog result, and only then closes the controller. If any of the
early steps fails, the dialog remains open.

\api{FormDialog.accept} implements this shape. It is an Event Dispatch Thread
operation because it inspects Swing editor state and mutates command state. It
returns a boolean rather than throwing for ordinary validation failure. A false
return means the current input did not produce a committed result. The caller
or validation binding can then focus the first problem, update a summary, or
leave the rejecting editor in place. A true return means the terminal
\api{DialogResult} is accepted and the controller is closed.

\begin{principlebox}{OK is conditional}
OK should be a request to produce a valid result. The form model decides
whether that result currently exists.
\end{principlebox}

The important invariant is one-way motion. Once OK succeeds, the controller has
reached a terminal closed state. The OK, Apply, Revert, and Cancel commands
become disabled. The accepted \api{DialogResult} is final. The shell may then
dispose the native peer, and the caller may continue with the result. This
one-shot rule makes tests and cleanup much simpler. There is no second OK after
a successful OK, and there is no hidden branch in which the controller closes
but the result remains mutable.

Apply is deliberately different. It uses the same active-editor and
committability path as OK, but it is non-terminal. On success it creates a
result snapshot, stores it as the last applied result, accepts the current form
values as the new dirty baseline, refreshes command state, and keeps the
controller open. Apply is useful when the user should see an effect while
continuing to edit: changing preview settings, applying filters, updating a
workspace option, or committing an intermediate page of a long configuration.
It is not a decorative sibling of OK.

The dirty baseline after Apply is the rule that must be stated explicitly.
Suppose the user changes three fields, presses Apply, changes one more field,
and then presses Cancel. In most desktop applications, the first three changes
remain because Apply committed them to the surrounding workflow; Cancel
abandons only changes since the last successful Apply. In a few specialized
workflows, Cancel may revert everything since the dialog opened. Both policies
can be defensible, but only one can be the user's mental model. Corusco's
default controller supports the conventional rule by calling
\api{acceptCurrentValues} after a successful Apply. After that point, user
Cancel returns the last applied value as an accepted result and abandons only
post-Apply edits. A program that chooses a different policy should not hide it
in a listener; it should name it and test it.

Revert is the separate operation for the other mental model: discard everything
since the dialog opened, including values already committed by Apply. It cannot
be implemented by ordinary form reset after Apply, because Apply deliberately
moved the reset baseline. Corusco therefore models Revert as an optional
\api{RevertPolicy}. The policy may restore an in-memory snapshot, call a
compensating save, or refuse because the surrounding workflow cannot honestly
undo the applied change. When it succeeds, the terminal \api{DialogResult} is
reverted rather than accepted, so callers do not confuse Apply--Cancel with
Apply--Revert.

Cancel is also a decision. It is not simply \api{dispose}. A clean dialog should
cancel quickly. A dirty dialog may need to ask whether user edits may be
discarded. \termdef{Dirty state}{for a dialog} is an aggregate answer to the
question ``does this activation contain user work that cancellation would lose?''
It may be derived from field dirty flags, table edits, nested editors, uploaded
files, or presenter state. Corusco represents that aggregate answer with
\api{DirtyState}. The controller asks it only for user cancellation, not for
lifecycle cleanup.

\termdef{Cancel confirmation}{in Corusco} is application policy invoked before
dirty edits are discarded. The \api{CancelConfirmation} hook may show a
\api{JOptionPane}, call an application dialog service, consult a presenter, or
return a canned decision in a test. It returns true to continue cancellation
and false to leave the controller open. The hook is intentionally not a Swing
component type. A library should not require every application to use the same
visual confirmation widget.

\begin{figure}[htbp]
\centering
\includegraphics[width=0.86\linewidth]{\dialogCancelFigure}
\caption{Dirty cancel confirmation is part of the modal protocol, not a private trick in the Cancel button listener.}
\label{fig:dialog-dirty-cancel}
\end{figure}

Figure~\ref{fig:dialog-dirty-cancel} shows the visible effect of this policy.
The user sees an ordinary question, but the controller sees a named sequence:
query dirty state, ask confirmation only when dirty, reset the form only after
confirmed cancellation, store a cancelled result, close, and disable commands.
That sequence should be shared by the Cancel button, the close-box branch if it
means user cancellation, and the Escape key. A codebase that implements dirty
confirmation three times will eventually have three slightly different
answers.

Programmatic close is different from user cancellation. \api{FormDialog.close}
forces controller cleanup as a cancelled dialog and bypasses dirty
confirmation. This is correct for owner-window disposal, application shutdown,
test cleanup, and lifecycle teardown. At that point the application is no
longer asking the user whether to discard a change; it is releasing resources
because the view lifetime has ended. Showing a dirty confirmation while the
owner window is already closing is one of those small UI defects that makes an
application feel as if its parts do not agree about who is in charge.

The controller also exposes OK, Apply, Revert, and Cancel as mutable commands. Commands
matter because a dialog operation is rarely invoked from exactly one button.
OK may be invoked by a button, a default-button Enter action, a toolbar action
in an embedded editor, or a test. Cancel may be invoked by a button, Escape, a
window close operation, or application navigation. If each surface owns its own
listener, command enablement and semantics drift. If each surface adapts the
same command, the operation has one identity, one enabled state, one resource
key, and one place to test.

\begin{lstlisting}[caption={Dialog buttons adapt dialog commands.}]
JButton ok = new JButton(new SwingActionAdapter(dialog.okCommand(), resources));
JButton apply = new JButton(new SwingActionAdapter(dialog.applyCommand(), resources));
JButton revert = new JButton(new SwingActionAdapter(dialog.revertCommand(), resources));
JButton cancel = new JButton(new SwingActionAdapter(dialog.cancelCommand(), resources));

buttonRow.add(ok, "tag ok");
buttonRow.add(apply, "tag apply");
buttonRow.add(revert, "tag other");
buttonRow.add(cancel, "tag cancel");
\end{lstlisting}

The resource keys behind the standard commands give the application a single
place for button text and other command metadata. The action keys give tests
and generated companions stable identities. The basic command enabled state
follows \api{refreshCommandState}: OK and Apply are enabled only while the
controller is open and the form is committable; Revert is enabled only while
the controller is open and the configured policy reports that restoration is
available; Cancel is enabled while the controller is open. If the form model
does not expose observable command-state changes, call \api{refreshCommandState}
after validation or parse state changes. This is a small explicit call, but it
is better than hidden listeners that revalidate the dialog at surprising times.

Button presentation state can be more precise than command availability. A
presenter may expose \api{ComponentStateModel} instances for Apply and Revert
through \api{FormDialogActionState}. Apply usually follows current-baseline
dirty state: it is enabled when there are unapplied changes and the form is
committable. Revert follows pre-dialog dirty state: it may remain enabled
after Apply, because the transaction still differs from the opening snapshot.
These action models are presentation state. They should not appear in the
accepted domain result.

There is one more rule that is easy to miss: \api{toResult} must return a real
result. A null result is not a cancelled result and not an invalid result; it is
a programming error. The form model already has \api{isCommittable} and
\api{problems} for ordinary invalid input. By requiring non-null committed
values, the controller avoids a three-valued outcome in which accepted-null and
cancelled are accidentally confused. If a domain result is naturally optional,
the accepted value should be an explicit optional domain object, not a null
reference smuggled through the dialog protocol.

The controller can be read as a small state machine. It starts open with a
cancelled result. OK may either leave it open or move it to accepted and closed.
Apply may either leave it unchanged or update the last-applied snapshot while
remaining open. Cancel may leave it open if dirty confirmation refuses, move it
to cancelled and closed when nothing has been applied, or move it to accepted
and closed with the last applied value. Revert may move it to reverted and
closed when the policy succeeds. Close always moves an open controller to
cancelled and closed without asking confirmation. Every transition refreshes
commands so the visible controls do not advertise actions that the controller
will ignore. The state machine is small enough to test exhaustively, and that
is one of the main reasons to keep it out of anonymous listeners.

The state machine also gives a precise meaning to idempotence. Calling Cancel
after the controller is already closed is harmless. Calling Close after it is
closed is harmless. Closing a keyboard binding or lifecycle twice is harmless.
Idempotence is not an excuse to call methods randomly; it is protection for
real shutdown paths in which the user closes a window, a parent disposes, and a
test finally block also tries to clean up. A modal workflow should end once,
but cleanup requests may arrive more than once.

Applications sometimes ask whether OK should be disabled while the form is
invalid or enabled and then explain problems after the user presses it. Both
policies can be reasonable. Disabling OK can reduce futile attempts, but it
must be paired with visible validation state; otherwise the user sees a dead
button and no explanation. Leaving OK enabled can be useful when validation is
expensive or when the user expects to press OK to discover what remains. The
controller supports command enablement from \api{isCommittable}; the
application chooses how eagerly the model updates that answer. The chapter's
examples use ordinary immediate validation, but the protocol does not require
every field to validate on every keystroke.

\begin{migrationbox}{From button listeners to commands}
Start by moving the existing OK logic into a method that returns a boolean.
Then make that method commit active editors, ask the form model for
committability, and create a result only on success. Next adapt the method as a
\api{FormDialog} command. Finally wire buttons, Enter, Escape, and tests to the
same command objects.
\end{migrationbox}

\section{The edit boundary}

\termdef{Active editor}{in Swing} is an editor component that currently holds a
candidate value for another component. The most familiar example is a
\api{JTable} cell editor. While the user edits a table cell, the editor
component may show the new text even though the table model still contains the
old value. \api{JFormattedTextField} and \api{JSpinner} have a related problem:
the visible text may not yet have been parsed and committed to the component's
value. A dialog that reads its model before flushing active editors can silently
lose the user's last edit.

The last-edit bug is common because the ordinary manual test path does not
always reveal it. A tester types a value, tabs out of the field, and presses
OK; the focus transfer commits the edit. A user types a value and immediately
presses OK; the active editor may still hold the candidate value. A tester
edits a table, clicks another row, and presses OK; the table stops editing. A
user edits a table cell and presses the dialog's default button; the model may
not see the change unless the dialog stops editing first. The difference is a
single gesture, and the data loss is hard to notice because the dialog closes.

\api{FormDialog} handles this boundary through \api{commitActiveEditor}. The
default implementation delegates to \api{SwingEditors.commitActiveEditor},
which inspects the focus owner inside the dialog root, commits focused
formatted fields and spinners, stops the focused table editor, and also walks
tables below the root to stop any active table editing. A false return means an
editor rejected the commit, so OK and Apply remain blocked and the controller
stays open.

\begin{pitfallbox}{The Last Edit Bug}
If a dialog loses only the final table-cell or formatted-field edit when OK is
pressed, active-editor commit is missing from the OK path.
\end{pitfallbox}

It is tempting to treat active-editor commit as merely a table issue. That is
too narrow. Any component that separates text editing from semantic value
commit can have this boundary. A formatted decimal field may display a string
that cannot be parsed; a date field may contain a partial date; a spinner may
hold a typed value that is not yet committed; a custom GIS coordinate editor
may have a temporary local point; a color editor may have a popup selection
that has not yet updated the model. The correct rule is general: before the
dialog creates a result, every editor under the root must either commit its
candidate value or reject the commit visibly.

Do not respond to a rejected active-editor commit by showing a generic error
dialog. The editor already knows why it rejected the value more precisely than
the dialog controller does. A formatted field can keep focus and show its
parse feedback. A table cell editor can remain active. A generated binding can
update a problem set. The controller's job is to stop the modal transaction
from completing while the candidate value is unresolved.

The form model's \api{isCommittable} check is the next boundary. It asks
whether the model, after editor commit, can produce a result. In a simple form
this may be equivalent to ``there are no error problems.'' In a richer form it
may also depend on asynchronous validation state, required selections,
cross-field consistency, or a pending lookup. A warning may still permit
commit; an error normally does not. The distinction should be encoded in the
problem model and in \api{isCommittable}, not recreated in each OK listener.

\termdef{Parsing}{in form workflows} is the conversion of user text into a
typed value. \termdef{Validation} is the evaluation of a typed value, or a
group of typed values, against rules. The two concerns often occur near each
other, but they should not be collapsed. A text field containing \api{abc} for
a decimal amount has a parse problem; there is no decimal amount yet. A parsed
decimal amount of \api{-5} may have a validation problem if the business rule
requires non-negative credit. A dialog can present both as problems, but the
code should retain the distinction because recovery differs. A parse problem
usually focuses the field and preserves text. A validation problem may involve
another field, a table row, or a form-wide invariant.

\begin{lstlisting}[caption={A minimal active-editor test shape.}]
JPanel root = new JPanel();
JTable table = new JTable(model);
table.setDefaultEditor(Object.class, new RejectingCellEditor());
root.add(table);
table.editCellAt(0, 0);

FormDialog<TestForm, String> dialog = new FormDialog<>(form, root);

assertThat(dialog.accept()).isFalse();
assertThat(dialog.isClosed()).isFalse();
assertThat(table.isEditing()).isTrue();
\end{lstlisting}

This test is small, but it guards an important user promise. If an editor
rejects a candidate value, the dialog must not pretend that a result exists.
The test does not click a native window. It constructs the relevant Swing
objects on the EDT, drives the controller, and asserts that the modal decision
did not complete. That is the right level for most dialog workflow tests.

Subclassing \api{FormDialog} is supported for narrow active-editor
customization. A screen with specialized editors may override
\api{commitActiveEditor} to integrate a custom editor framework or to improve
test control. The override should preserve the controller invariants: remain
EDT-confined, do not mutate result state, do not close the dialog, and return
false when the current editor cannot publish a value. Once the editor boundary
is crossed successfully, normal form committability and result creation should
continue through the standard controller path.

This is also where generated code should be modest. A generated companion can
know field keys, component keys, parse policies, validation metadata, and
binding installation. It should not hide a dialog's active-editor policy in
unreviewed generated listener code. The dialog controller gives generated and
handwritten parts a common boundary: generated bindings make the model honest;
the controller asks the whole root to settle before accepting.

There is an important performance angle in large dialogs. Committing active
editors should be precise, not a full reparse of every field in a five-tab
editor. The focused editor and actively editing tables need attention because
they may contain unpublished candidates. The form model may then validate the
affected fields or cross-field rules according to its own dependency graph.
Running every remote validator, rebuilding every table model, and repainting
the whole dialog on each OK attempt makes the dialog feel arbitrary. Users
accept a short pause when they know a save is happening; they do not accept a
pause merely because the program could not distinguish the active edit from the
rest of the screen.

Table dialogs deserve special suspicion. A table has model coordinates, view
coordinates, selection state, sorting, filtering, renderer state, editor state,
and sometimes row-level validation. The accepted result should not be assembled
from selected view rows without translation. The focus resolver should not
scroll to a model row by treating it as a view row. The active-editor commit
path should stop editing before the model is read. A small table in a dialog is
still a \api{JTable}; the scale may be smaller, but the coordinate systems are
the same as in a full table screen.

Custom components should expose an honest commit contract. If a component can
hold a candidate value that is not yet in the model, provide a method or binding
that commits or rejects it. If it cannot reject, the method can always return
true. If it can reject, the component should preserve the candidate and present
the reason. Do not make the dialog guess by reading private editor state. The
dialog root is the traversal boundary; the component remains the authority over
its own local edit semantics.

\section{Validation that guides recovery}

\termdef{Problem set}{in Corusco} is a structured collection of problems,
ordered and filtered by severity, target, and code. A problem is more than a
localized message. It has a stable problem code, a severity, a target, and
text suitable for the user. A dialog that stores only the string
\api{"Name is required"} has already thrown away information that focus
recovery, tests, help, analytics, and generated companions could use.

\termdef{Problem target}{in this context} is the structured address of the
thing that needs attention. A target may be the whole form, a typed field key,
a typed component key, a table row and column, or another application-defined
location. The target is not necessarily a Swing component. This distinction
lets the same validation rule exist in core code while different Swing screens
map the target to their own components.

\api{FormDialogValidationBinding} is the Swing presentation companion to the
controller's commit checks. It reads the current form problems, writes a compact
summary into a \api{JLabel}, and can ask a \api{ProblemFocusResolver} for the
component that should receive focus. It does not run validation. It does not
create dialog results. It does not decide whether OK is enabled. It presents
the current problem state at one place in the dialog.

\begin{lstlisting}[caption={Validation summary and focus binding.}]
FormDialogValidationBinding validation =
        FormDialogValidationBinding.install(
                dialog,
                summaryLabel,
                ProblemFocusResolver.componentTargets(components));

validation.refresh();
if (!dialog.accept()) {
    validation.focusFirstProblem();
}
\end{lstlisting}

The binding is intentionally pull-based: call \api{refresh} after form state
changes or after a generated presenter runs validation. In a future form model
with observable problem state, a listener could call the same refresh method.
The important architectural rule is unchanged: the form model owns the
problems; the validation binding owns one Swing presentation of them. A tooltip
binding, a red border behavior, a status line, a table marker, and a dialog
summary can all read the same structured data.

The summary policy is compact by design. When there are no problems, the label
is blank. When there is one problem, the label displays that message. When
there are several, it displays a count plus the most severe message. This is
not meant to be the only possible presentation. It is a sane default for the
bottom of a dialog. A larger screen might also show a problem list, group
problems by tab, or mark table rows. Corusco's contribution is that all such
presentations can agree on the same severity ordering and targets.

Focusing the first problem is harder than it looks. The first severe problem
may be form-wide and have no component. A field problem may live in a tab that
is currently hidden. A table-cell problem may require translating model
coordinates to view coordinates after sorting, scrolling the row into view,
selecting a cell, and starting an editor. A problem in a generated subform may
belong to a component that is created lazily. The validator should not know any
of that. Validators describe what is wrong; focus resolvers know how this
particular Swing screen can lead the user to the place where the problem can be
fixed.

\api{ProblemFocusResolver.componentTargets} covers the simple and common case:
component-targeted problems map to typed component keys, and the resolver
returns the corresponding Swing component. Generated code can produce that map
from component keys. Handwritten code can build it explicitly. No reflection,
localized labels, JavaBeans property names, or table view indexes are required.
For more complex screens, write a resolver that performs the needed navigation
before returning or focusing a component.

\begin{lstlisting}[caption={A resolver for a tabbed form and a table cell.}]
ProblemFocusResolver resolver = problem -> {
    if (problem.target().equals(CustomerTargets.billingAddress())) {
        tabs.setSelectedComponent(billingPanel);
        return Optional.of(streetField);
    }
    if (problem.target() instanceof CustomerTargets.TableCell cell) {
        int viewRow = table.convertRowIndexToView(cell.modelRow());
        int viewColumn = table.convertColumnIndexToView(cell.modelColumn());
        table.changeSelection(viewRow, viewColumn, false, false);
        table.scrollRectToVisible(table.getCellRect(viewRow, viewColumn, true));
        return Optional.of(table);
    }
    return Optional.empty();
};
\end{lstlisting}

The example is deliberately explicit. A focus resolver is a navigation adapter,
not a second validation engine. It should not inspect field values and decide
what is wrong. It should translate problem targets into user recovery. Tests
for this code should assert that the right tab is selected, the right table
cell is visible, and the right component receives a focus request.

For multi-form dialogs, each child view can expose its own explicit resolver.
The parent presenter composes those resolvers in the same order as the child
forms and wraps each child resolver with the preparation needed to reveal that
child. A problem in a hidden tab should first select the tab. A problem in a
collapsed section should first expand the section. The validator still reports
a typed target; the presenter owns the navigation.

\begin{lstlisting}[caption={Reveal the child view before focusing its problem.}]
ProblemFocusResolver resolver = ProblemFocusResolver.firstOf(
        ProblemFocusResolver.withPreparation(
                () -> tabs.setSelectedComponent(identityPanel),
                ProblemFocusResolver.fieldTargets(Map.of(
                        CustomerIdentityFields.NAME.asFieldKey(),
                        identityView.nameField()))),
        ProblemFocusResolver.withPreparation(
                () -> tabs.setSelectedComponent(securityPanel),
                ProblemFocusResolver.fieldTargets(Map.of(
                        SecuritySettingsFields.PASSWORD.asFieldKey(),
                        securityView.passwordField()))));
\end{lstlisting}

This shape keeps the ownership boundaries clean. The core problem target is a
typed field key. The child view supplies a component. The parent presenter
selects the child surface. The validation binding only requests focus on the
component it was given.

Validation failure and unexpected failure deserve different presentations.
Validation failure is ordinary user input. It belongs in problem summaries,
field decorations, tooltips, and focus recovery. Unexpected failure during a
save or apply operation belongs to workflow error handling: status text, retry,
diagnostics, or an application error dialog. A repository exception is not a
field validation problem unless the field can help the user fix it. Conversely,
a missing required field is not an exceptional failure and should not produce a
stack trace dialog.

Asynchronous validation needs another layer of care. Consider a dialog that
checks whether a customer code is unique. The user types \api{ACME}, the
program starts a background lookup, the user changes the code to \api{ACME-2},
and the first lookup returns late. If the old result writes a problem into the
current form, the dialog lies. The usual solution is a generation counter or a
request token: each submitted validation captures the input version, and only
the latest matching generation may publish problems. The dialog's OK command
should also know whether a required asynchronous validation is pending. A
pending required validation may disable OK or make \api{isCommittable} false
until the result arrives.

Corusco's task and value facilities help here because they give pending state,
busy state, cancellation, and problem publication explicit owners. The dialog
controller still should not run the uniqueness query. It should ask the form
model whether the current state is committable. The presenter or validation
service may update the form's problem set and call \api{refreshCommandState}
and \api{FormDialogValidationBinding.refresh}. Keeping these roles separate
prevents a common bug in which the OK listener both starts a validation lookup
and closes the window on a stale answer.

Problem codes are also a documentation device. A localized message may change
from ``Name is required'' to ``Enter a customer name'' without changing the
meaning of the rule. Tests that assert message text become brittle; tests that
assert problem code and target remain stable. Help systems can link a code to a
longer explanation. Analytics can count how often a particular validation rule
blocks a workflow without parsing user text. Generated documentation can list
the validation rules for a form. These advantages disappear when the problem is
only a string assigned to a label.

Severity needs the same discipline. An error normally blocks OK. A warning may
permit OK while still asking the user to review the field. An informational
problem may explain a derived value or a policy. A dialog summary that treats
all severities as errors is noisy; a controller that treats warnings as
non-committable is paternalistic unless the business rule truly requires it.
The form model is the place to encode which severities block commit. The
presentation is the place to help the user understand them.

Multiple presentations of the same problem state should not compete. A field
border, a tooltip, a summary label, and a table marker can all exist together,
but they should be consistent. If the border says the value is invalid and the
summary says the form is valid, the user will trust neither. Corusco's binding
and behavior scopes are intentionally small so a screen can install several
presentations that all read the same \api{ProblemSet}. The result is richer UI
with less duplicated validation code.

\begin{coruscobox}{Structured Validation State}
Corusco treats validation as structured state rather than scattered strings.
The dialog controller asks whether the form is committable; validation bindings
present problem state; focus resolvers translate problem targets into Swing
recovery. The same problems can drive tooltips, borders, summaries, disabled
commands, and tests.
\end{coruscobox}

\section{Keyboard, focus, and the native shell}

Swing keyboard behavior is powerful because it is component-local. The same
power makes dialogs easy to wire inconsistently. A root pane can have a default
button. A text field can use Enter to fire an action. A text area can use Enter
to insert a newline. A combo box can use Enter to accept a popup selection. A
table can use Enter to finish editing or move selection. Escape may close a
popup, cancel an editor, or cancel the dialog. A correct dialog should respect
local component meaning before turning a key into a modal decision.

\api{FormDialogKeyboardBinding} installs the shared root-pane part of the
policy. It maps Escape in the root pane's \api{WHEN\_IN\_FOCUSED\_WINDOW} input
map to the dialog's cancel command and can install a default OK button. The
binding does not own cancellation semantics. It routes keyboard events to the
same command used by visible buttons. Therefore Escape respects dirty
confirmation, command enablement, one-shot close, and tests in exactly the same
way as the Cancel button.

The binding also records what it changed and restores it on close. It remembers
the previous Escape input-map value, the previous action-map entry, and the
previous default button. Closing the binding restores those values, and closing
twice is harmless. This may sound like housekeeping, but it is essential when
dialogs are embedded in reusable shells, internal frames, or test harnesses.
Root panes are shared infrastructure; a modal activation should not leave its
Escape mapping behind.

\begin{lstlisting}[caption={Keyboard binding belongs to the dialog lifecycle.}]
FormDialogLifecycle lifecycle = FormDialogLifecycle.of(dialog);
lifecycle.addBinding(FormDialogKeyboardBinding.install(
        rootPane,
        dialog,
        okButton));
\end{lstlisting}

Default-button behavior deserves review rather than superstition. In many
simple dialogs, Enter should invoke OK from ordinary single-line fields. In a
multi-line text area, Enter should insert a newline. In a table cell editor,
Enter may stop editing or move to another row. In an editable combo box, Enter
may accept the current popup item. Swing's input maps and component actions
usually handle these local meanings before the root-pane default button becomes
relevant, but custom components can surprise you. The rule is not ``Enter
always accepts''; the rule is ``the accepted Enter path must eventually invoke
the same OK command as the button.''

Focus after validation failure is similarly a dialog responsibility but not a
dialog-controller responsibility. The controller knows that OK failed. The
validation binding knows the current problems. The focus resolver knows how to
map problem targets to components. The native shell knows which window is
active. These objects cooperate. A single monolithic OK listener often tries to
do all of this: validate, set label text, select tabs, request focus, decide
dirty cancel, and close. Such a listener cannot be reused and is tedious to
test. More importantly, it tempts each dialog author to invent a slightly
different focus policy.

The native \api{JDialog} shell should remain small. It sets the owner, title,
modality, content pane, default close operation, packing, placement, and native
disposal. It may adapt window-closing events to user cancellation or lifecycle
close according to application policy. It should not store the accepted domain
object in a public field. It should not rerun validation in a window listener.
It should not duplicate OK semantics for the close box. If the shell needs to
know whether the controller has accepted or cancelled, it can inspect the
controller result after the modal show returns.

There is a subtle distinction between hiding and disposing a native dialog.
Hiding a reusable dialog keeps component instances, listeners, and model
references alive. Disposing releases the native peer, but Java object
references still live if listeners or owners retain them. For most modern Swing
business applications, one controller and one view activation per modal
interaction is simpler than reusing the same dialog instance many times. Create
the form model, create the panel, install bindings, show, collect the result,
close lifecycle, and let the objects become unreachable. Reuse should be a
measured optimization, not the default.

Owner relationships affect modality and focus. A dialog with no owner can hide
behind the main window or keep the wrong window blocked. A dialog owned by a
stale window can behave strangely during shutdown. A dialog opened from a table
row, a secondary tool window, or a document tab should normally be owned by the
nearest active window in that workflow. Corusco does not choose the owner for
you because application window topology is policy. It does make the controller
small enough that shell ownership can be reviewed separately from form
semantics.

Accessibility also starts here. Buttons adapted from commands should have
stable text and enabled state. Fields should have labels and accessible names.
Validation summaries should not be the only accessible presentation of an
error; field-specific descriptions or tooltips may be needed. Focus recovery
should lead keyboard and assistive-technology users to the same place a mouse
user would inspect visually. A dialog that paints a red border but gives no
focus path or accessible description has solved only part of the problem.

Window-closing policy should be written down. Some applications treat the close
box exactly like Cancel, including dirty confirmation. Others treat owner
disposal or application shutdown as lifecycle close and bypass confirmation.
The difference depends on who initiated the close. If the user clicks the close
box on the dialog, Cancel semantics are usually right. If the owner window is
being disposed because the application is shutting down, lifecycle cleanup is
usually right. A shell helper can make this distinction once. A collection of
anonymous window listeners usually cannot.

Focus activation has its own timing. Requesting focus before a native dialog is
shown often fails because there is no active window yet. Requesting focus after
validation failure usually succeeds only if the target component is visible,
enabled, displayable, and inside the active window. A focus resolver can select
tabs and scroll tables, but the shell may still need a later focus request on
initial show. Keep these two cases separate: initial focus is part of shell
presentation; validation focus is part of problem recovery.

The keyboard binding is deliberately conservative. It does not try to rewrite
every component's input map. Rich components often have legitimate local key
bindings, and overwriting them from a dialog helper would create harder bugs
than it solves. The root-pane binding establishes the shared modal route; custom
components remain responsible for local editing behavior. This is the same
architecture as painting and layout throughout Swing: local components keep
their contracts, and shared containers coordinate only the surrounding policy.

\begin{pitfallbox}{Two Cancel Paths}
If Escape discards edits immediately but the Cancel button asks for
confirmation, the dialog has two cancellation paths. Route both through the same
cancel command.
\end{pitfallbox}

\section{Dialog bodies with MigLayout}

A dialog body is a design object. It expresses which facts belong together,
which fields grow, which labels align, where validation feedback appears, and
which controls are visually secondary. For ordinary Swing business dialogs,
MigLayout is the right default because it describes this geometry directly.
The source reads as a form: label, field, label, field, summary, buttons. The
alternative in old Swing code is often a nest of \api{BorderLayout},
\api{BoxLayout}, and \api{GridBagLayout} panels whose geometry is distributed
across too many objects to review easily.

\begin{lstlisting}[caption={Typical MigLayout form body.}]
JPanel body = new JPanel(new MigLayout(
        "insets dialog, fillx",
        "[][grow, fill]",
        ""));

body.add(new JLabel("Name:"), "right");
body.add(nameField, "growx, wrap");
body.add(new JLabel("Credit limit:"), "right");
body.add(creditLimitField, "growx, wrap");
body.add(summaryLabel, "span 2, growx, wrap");
body.add(buttonBar, "span 2, right");
\end{lstlisting}

The layout constraints are not decoration. \api{insets dialog} says the panel
is a dialog body and should use dialog-like margins. \api{fillx} says the panel
may use horizontal space. The column constraint \api{[][grow, fill]} says
labels have natural width while fields receive extra width and fill their
cells. \api{wrap} says the next component starts a new row. \api{span 2} says
the summary and button row belong across both columns. These words are compact,
but they are also a readable geometry language.

Button rows are one place where MigLayout's platform vocabulary is particularly
useful. Tags such as \api{tag ok}, \api{tag cancel}, and \api{tag apply} let the
layout manager arrange buttons according to platform rules where supported.
Even when an application chooses a fixed order, the tags document semantic
roles. The buttons themselves should still be command adapters; the layout tag
does not own behavior.

A validation summary should usually be close to the action row, because it
explains why the requested action did not complete. Field-level decorations
should remain near fields. Form-wide problems may appear in the summary, a
problem list, or a banner. Avoid modal error popups for ordinary validation.
They interrupt the edit flow and force the user to dismiss a second dialog
before returning to the first. A validation summary plus focus recovery is
quieter and more useful.

Growing components need explicit thought. A single-line text field may grow
horizontally but not vertically. A table, list, text area, preview, or map
component may need both horizontal and vertical growth. A dialog with no
vertical grow target may look cramped at large sizes and may waste space. A
dialog with every row growing may look unstable. MigLayout lets the source say
which rows grow and which remain at preferred height. For a generated form,
the generator can provide field metadata, but the screen author should still
decide growth.

Do not let generated code own the whole layout merely because it can. Generated
companions are excellent for field keys, resource keys, model/view contracts,
binding installation, validation metadata, component keys, and tests. Layout is
often where product decisions live: grouping, density, ordering, progressive
disclosure, and visual rhythm. A generated layout may be acceptable for an
administrative prototype, but a durable application screen deserves handwritten
MigLayout or a generated layout that remains easy to inspect and override.

The same caution applies to resource text. Labels, button text, validation
messages, tooltips, and help topics should come from resources or generated
keys where appropriate, but the body should still show the relationship between
the label and the component. A panel whose source only says
\api{GeneratedCustomerForm.create()} may compile, but it hides the screen from
review. A panel that says \api{body.add(view.nameLabel())} followed by
\api{body.add(view.nameField())} keeps the design visible while still using
generated components.

Dialog density should match workflow frequency. A rare destructive
confirmation can be sparse. A daily operational editor should be compact,
aligned, and predictable. Large business dialogs often need sections, but too
many visual boxes make the dialog heavy. Use full-width separators or labeled
groups sparingly. A tabbed dialog should have tabs only when the user benefits
from separate work areas; tabs are not a substitute for simplifying the model.
If validation can target hidden tabs, the focus resolver must be able to select
them.

The best layout test is not a pixel-perfect assertion. It is a small set of
deterministic screenshots and smoke tests. Render the dialog body at the sizes
the application supports. Verify that labels do not wrap unexpectedly, buttons
remain visible, summaries have space, table rows are not crushed, and the first
invalid field can receive focus. The companion screenshot harness uses FlatLaf,
fixed dimensions, deterministic data, and generated PNGs precisely so the book
can show realistic Swing surfaces instead of pseudocode.

A dialog body should also have a resize story. Fixed-size dialogs are sometimes
appropriate for tiny confirmations, but editors with tables, lists, previews,
or text areas should usually resize. The resize story says which component
absorbs width, which component absorbs height, and which parts remain at
preferred size. MigLayout makes this story explicit with grow and fill
constraints. If the story is missing, the dialog may look fine at the
developer's scale factor and poor at 125 percent, 150 percent, translated text,
or a different look and feel.

Localization tests belong to layout tests. German labels, Czech messages, and
long generated resource names can reveal fragile columns. A validation summary
with a count and a long message may need to wrap or elide. Buttons may become
wider. A design that relies on exact English string lengths is not a design;
it is an accident. Because Corusco encourages stable resource keys, the visual
text can change without changing the layout code, but the layout still must be
tested with realistic strings.

Example references should be intentionally short. A reader should see
\dialogExample{} or \dialogValidationExample{}, not a full package path that
wraps across the page. The repository structure can move; the class name
remains a readable entry point. Use the same convention for new examples: name
the example by the didactic program, and let the examples index or build system
point to its physical location.

\begin{examplebox}{Dialogs}
\dialogValidationExample{} and \dialogActiveEditorExample{} are the compact
examples behind the dialog figures. They use ordinary Swing components and
short Corusco controllers so that chapter listings can stay small.
\dialogApplyRevertExample{} is the focused behavior example for Apply-Cancel
and Revert baselines.
\end{examplebox}

\section{Composite form sessions}

Large dialogs often contain several independently meaningful forms: identity,
billing, security, contacts, advanced options, or a generated form next to a
handwritten table editor. Flattening every field into one generated form makes
those screens harder to own and test. Corusco handles this by keeping the
dialog controller unchanged and composing the semantic form model instead.

\api{AbstractCompositeFormModel} is the core-side base for this pattern. It is
still a \api{FormModel}; it simply owns an ordered list of child form models.
Problem aggregation follows child order and then parent validation. Reset and
baseline acceptance flow to every child. Result creation remains an explicit
parent operation.

\begin{lstlisting}[caption={A handwritten parent session over generated children.}]
final class CustomerSettingsSession
        extends AbstractCompositeFormModel<CustomerSettingsResult> {

    private final CustomerIdentityFormModel identity;
    private final SecuritySettingsFormModel security;

    CustomerSettingsSession(
            CustomerIdentityFormModel identity,
            SecuritySettingsFormModel security) {
        super(identity, security);
        this.identity = identity;
        this.security = security;
    }

    @Override
    protected CustomerSettingsResult createResult() {
        return new CustomerSettingsResult(
                identity.toResult(),
                security.toResult());
    }
}
\end{lstlisting}

Parent validation is the reason the parent exists. It can express rules such
as ``external authentication is not allowed for inactive customers'' without
making the security form know about the customer form. Parent validation
should inspect child field state or child committability. It should not
materialize child results unless the child is already committable; the
assembler is the normal place for child \api{toResult} calls.

The existing \api{FormDialog} can host the composite directly because the
composite is still a form model:

\begin{lstlisting}[caption={The dialog sees one form model.}]
DirtyState dirty = DirtyStates.any(
        () -> identity.name.isDirty(),
        () -> security.password.isDirty());

FormDialog<CustomerSettingsSession, CustomerSettingsResult> dialog =
        new FormDialog<>(session, root, dirty, confirmation);
\end{lstlisting}

The presenter side should compose generated child presentation models in the
same spirit. A parent presenter can keep \api{CustomerIdentityPresentationModel}
and \api{SecuritySettingsPresentationModel} beside dialog-level action state,
tab selection, and problem routing. None of that presentation state belongs in
the committed result. Swing-specific generated companions remain optional:
packages that opt into \api{SwingCompanionPackage} can use generated views and
bindings, while the core session only depends on core form and presentation
models.

\begin{examplebox}{Composite dialogs}
\multiFormDialogExample{} demonstrates a generated-child, handwritten-parent
dialog session. It composes generated form and presentation models, derives
Apply and Revert state from separate dirty policies, routes child field
problems through typed focus resolvers, and keeps generated Swing bindings
optional.
\end{examplebox}

\section{Lifecycle and background work}

\termdef{Dialog lifecycle}{in Corusco Swing} is the owned set of temporary
resources that live for one dialog activation. It includes field bindings,
validation bindings, keyboard bindings, behavior scopes, listener handles, task
services, detachable presenter state, and the dialog controller itself. The
lifecycle is not the native window and not the form model. It is the cleanup
owner beside them.

\api{FormDialogLifecycle} collects these resources. It is EDT-confined, like
the Swing objects it owns. A presenter creates the controller, creates the
lifecycle, registers installed bindings and services, and closes the lifecycle
when the dialog activation ends. The lifecycle closes registered resources in
reverse registration order, then closes the dialog controller in a finally
block. If a late resource is registered after the lifecycle is already closed,
the underlying resource scope closes or detaches it immediately. This rule
removes a surprising number of edge-case branches from presenters.

\begin{figure}[htbp]
\centering
\includegraphics[width=0.92\linewidth]{\dialogLifecycleFigure}
\caption{A dialog activation should own temporary resources explicitly. Cleanup closes bindings and services before the controller resets the form.}
\label{fig:dialog-lifecycle}
\end{figure}

The reverse order is not accidental. If a field binding is still installed
while the controller resets the form during cancellation, the binding may
publish reset values back into components or fire listeners after the view is
already considered dead. Closing bindings first removes the component-model
relationships. Closing the controller last records the cancelled result and
resets the model as part of terminal cleanup. If a binding cleanup throws, the
controller still closes, because a failed listener removal should not leave the
modal protocol open.

\begin{lstlisting}[caption={One lifecycle for a dialog activation.}]
FormDialog<CustomerForm, Customer> dialog = new FormDialog<>(form, body);
FormDialogLifecycle lifecycle = FormDialogLifecycle.of(dialog);

lifecycle.addBinding(fieldBindings);
lifecycle.addBinding(FormDialogValidationBinding.install(dialog, summary, resolver));
lifecycle.addBinding(FormDialogKeyboardBinding.install(rootPane, dialog, okButton));
lifecycle.addTaskService(taskService);
\end{lstlisting}

Background work inside dialogs is common: loading choices, validating a remote
identifier, previewing a report, applying a settings change, or saving a large
object. The dialog should remain responsive. Long work must not block the EDT.
At the same time, results must return to the EDT before they mutate Swing
components or form models. This boundary is the same as in full screens, but
dialogs make the lifetime sharper. If the user cancels the dialog while a task
is running, the task must not update a disposed view.

A task service registered with the dialog lifecycle gives the work a lifetime.
Closing the lifecycle closes the service according to the service's own
semantics. A generation counter or task handle should prevent stale results
from publishing into a newer activation. Busy state can disable OK or Apply,
show a progress bar, update status text, and install a busy overlay. The
controller should still ask the form model whether it is committable; pending
task state can contribute to that answer through the model or presenter.

Virtual threads in modern Java are useful here, but they do not remove the EDT.
A virtual thread can perform blocking IO for a validation lookup or save
operation, but it must hand results back to the EDT before touching Swing. A
structured task scope can organize related background work, but a dialog
activation still needs cancellation and stale-result suppression. The modern
feature changes the worker side of the boundary; it does not make Swing
thread-safe.

Busy UI should be proportional. A small lookup may need only a disabled OK
button and a status label. A longer operation may need a progress bar and a
Cancel current task button. A modal overlay can be useful when the whole form
is temporarily unavailable, but it should not prevent users from cancelling if
cancellation is safe. The busy chapter discusses these patterns in depth; the
dialog-specific rule is that busy state belongs to the same activation
lifecycle as the task that produces it.

Dialogs that perform a save on OK need special care. One reasonable pattern is
to let OK produce a result and close the dialog, then let the caller perform
the save in the surrounding screen with its own busy/error presentation.
Another pattern is to keep the dialog open while Apply or Save performs a
background operation, then close only after success. The second pattern is
valid when the user needs immediate feedback inside the modal context, but the
dialog must then distinguish validation failure from save failure and must
avoid accepting stale task results after cancellation. Do not mix the two
patterns accidentally.

Cancellation should travel down to background work whenever possible. A user
who cancels a dialog that is loading choices does not expect the loading task
to keep a database cursor open or to publish into a hidden combo box later.
Some work cannot be interrupted immediately, but the UI ownership should still
be cancelled: callbacks check whether the lifecycle is closed, task handles
are asked to cancel, and stale generations are ignored. This turns cancellation
from a visual event into a resource boundary.

The lifecycle is also the right place for temporary caches. A dialog may cache
lookup values, rendered previews, remote validation responses, or table filter
state for the duration of one activation. Keeping such caches in static fields
or long-lived presenters often creates stale state in the next activation.
Registering a detachable cache with \api{FormDialogLifecycle} states that the
cache is useful only while this modal edit exists. When the lifecycle closes,
the cache detaches and the next activation starts with a clean boundary.

Pay attention to cleanup side effects. Closing a binding should restore the
component state it owned and remove listeners it installed. It should not run
validation, save data, or show UI. Closing a task service should stop accepting
new work and arrange cancellation according to its contract. Closing the
controller should reset the form for cancelled cleanup. These operations are
ordered, but they should remain individually simple. Complex cleanup is a sign
that the dialog activation has been allowed to own too many unrelated things.

Resource ownership also includes generated behavior scopes. A generated form
binding plan may install value bindings, validation border behaviors, composed
tooltips, help hooks, and command adapters. These bindings are not global.
Register them with the dialog lifecycle or a behavior scope that is itself
registered. When the dialog closes, the scope should restore component state
that it changed, such as borders, tooltips, input maps, and listener lists. A
dialog should leave the application as if its temporary view had never existed.

\begin{coruscobox}{Modal Cleanup as a First-class Object}
Corusco makes modal cleanup a first-class object. \api{FormDialogLifecycle}
lets presenters register bindings, disposables, detachables, and task services
beside one controller, then close the whole activation in a predictable order.
\end{coruscobox}

\section{Plain Swing and the migration path}

A hand-written Swing dialog often begins with two button listeners and a field
named \api{result}. This is not foolish; it is the shortest path to a working
demo. The problem is that the pattern does not scale with the number of rules a
real dialog needs. The OK listener grows active-editor commit, validation,
problem summary, focus selection, result creation, exception handling, dirty
baseline updates, command enablement, and disposal. The Cancel listener grows
dirty confirmation and reset. The window listener grows another cancellation
branch. The root pane grows keyboard actions. Later a test discovers that the
Escape key does not match the Cancel button.

\begin{lstlisting}[caption={Typical handwritten OK/Cancel drift.}]
okButton.addActionListener(event -> {
    if (!form.isCommittable()) {
        summaryLabel.setText(form.problems().bySeverityDescending()
                .getFirst().message());
        return;
    }
    result = form.toResult();
    dialog.dispose();
});

cancelButton.addActionListener(event -> dialog.dispose());
\end{lstlisting}

This code omits active-editor commit, dirty confirmation, command enablement,
keyboard consistency, validation focus, lifecycle cleanup, Apply, and result
typing. Each omission is easy to fix in isolation. The damage appears when
every dialog fixes a different subset. Once that happens, the team no longer
has a dialog policy; it has dialog folklore.

The migration path is incremental. First move durable edit state out of Swing
widgets and into a form model. The existing panel can remain. Next make OK ask
the form model for committability and result creation instead of reading fields
directly. Then introduce \api{FormDialog} and adapt the existing buttons to
its commands. Add active-editor commit by giving the controller the real root
component. Add dirty confirmation only after the clean baseline is defined.
Move the validation summary and focus recovery into
\api{FormDialogValidationBinding}. Install root-pane keyboard behavior with
\api{FormDialogKeyboardBinding}. Register everything in
\api{FormDialogLifecycle}.

\begin{migrationbox}{A practical sequence}
Do not rewrite an old complex dialog in one step. Extract the form model, then
the controller, then validation presentation, then keyboard binding, then
lifecycle cleanup. After each step, delete one private flag, one duplicated
listener branch, or one manually synchronized button state.
\end{migrationbox}

This path works because Corusco does not require the layout to change. An old
dialog can keep its current \api{JPanel}. A screen with a carefully tuned
MigLayout body can keep it. A generated form can use generated field accessors
while the body remains handwritten. The controller only needs a form model, a
root component, and optional dirty/cancel hooks. This is important in mature
Swing applications, where a modal editor may contain domain-specific custom
components and years of usability adjustments.

The first useful test usually targets the controller without a native shell.
Construct a small form model on the EDT, create a root panel, create a
\api{FormDialog}, drive \api{accept}, \api{apply}, and \api{cancel}, and assert
the result, closed state, form reset calls, baseline calls, and command
enablement. This test proves the modal protocol. It does not depend on window
placement, native focus, or user input timing.

The second useful test targets validation presentation. Give the form model a
problem set, install \api{FormDialogValidationBinding}, refresh it, and assert
that the summary label contains the expected message. Provide a resolver that
maps a typed component target to a test component whose \api{requestFocusInWindow}
method records calls. Then assert that focusing the first problem chooses the
right component. This is much more direct than clicking through a modal window
and hoping focus is observable.

The third useful test targets active-editor commit. Use a table with a custom
cell editor that rejects \api{stopCellEditing}, or override
\api{commitActiveEditor} in a test subclass. Assert that OK and Apply return
false, the controller remains open, and no result is created. This test guards
the most common data-loss edge.

The fourth useful test targets keyboard and lifecycle. For keyboard behavior,
install \api{FormDialogKeyboardBinding} on a \api{JRootPane}, trigger the
Escape action from the input/action maps, and assert that the controller's
cancel command ran and dirty confirmation was honored. For lifecycle behavior,
register tracked bindings or disposables and assert reverse cleanup order and
idempotent close. These tests are not glamorous, but they catch the defects
that make dialogs feel unreliable.

Generated companions can add stronger tests. A generated form can expose typed
field keys, component keys, validation problem codes, and command keys. A test
can assert that each required field has a component, each validation problem
has a focus target, each generated command is adapted exactly once, and each
installed behavior is closed with the lifecycle. This is where generation is
most valuable: not by hiding Swing, but by giving the repeated facts stable
names that tests can use.

For legacy migration, an adapter layer is often better than direct surgery.
Wrap the old dialog's domain object in a form model that initially delegates to
the existing getters and setters, then gradually move state into Corusco
values. Replace one listener branch at a time. Keep the visual panel stable
while the modal protocol becomes explicit. This reduces the chance that a
behavioral migration is blamed on layout changes or resource text changes. Once
the protocol is tested, visual cleanup becomes less risky.

Do not migrate only the successful path. It is easy to replace OK with a
controller and leave Escape, close-box, Apply, validation focus, and lifecycle
cleanup untouched. That produces a dialog with a modern center and legacy
edges. Users usually find the edges. A migration review should exercise every
way the dialog can end, every way it can fail to end, and every temporary
resource it installs. The controller is the center, but the protocol is the
whole shape.

Screenshot-oriented examples should be compiled examples, not drawing-only
mockups. The figures in this chapter are deterministic panels, but they are
generated by the same example project that contains runnable dialog scenarios
and tests. This keeps the book honest. If a chapter describes active-editor
commit, there should be a test or example that drives active-editor commit. If
it describes dirty cancel, there should be a test or example that vetoes and
confirms cancellation. The screenshot is evidence of presentation; the tests
are evidence of behavior.

\begin{examplebox}{Dialogs}
\dialogExample{} demonstrates base OK, Apply, Apply-Cancel, and Cancel controller semantics.
\dialogApplyRevertExample{} demonstrates Apply-Cancel and optional Revert.
\dialogDirtyCancelExample{} covers dirty confirmation. \dialogKeyboardExample{}
covers Escape and default-button wiring. \dialogLifecycleExample{} covers
cleanup ownership. The book names the examples by class so that references stay
readable in print.
\end{examplebox}

\section{Worked review of a customer editor}

Consider a customer editor with name, credit limit, active status, billing
address, a table of contacts, and a server-side uniqueness check for the
customer code. This is large enough to expose the real problems, but small
enough that the modal protocol remains visible. The plain Swing version often
starts as a \api{JDialog} subclass with fields, table model, button listeners,
and a public \api{getCustomer} method. The Corusco version starts with a
question: what is the accepted result?

The result might be a record:

\begin{lstlisting}[caption={A dialog result should be explicit.}]
record CustomerEdit(
        CustomerId id,
        String name,
        BigDecimal creditLimit,
        boolean active,
        Address billingAddress,
        List<ContactEdit> contacts) {
}
\end{lstlisting}

The form model then owns editable values and problems for those facts. The
name field has raw text, parsed value if conversion is needed, dirty state, and
required validation. The credit limit field distinguishes parse failure from
negative-value validation. The active checkbox is simple but still participates
in dirty state. The contacts table has its own table model and perhaps
row-level problems. The uniqueness check publishes a pending state and then a
problem if the current code is already used. \api{toResult} creates the record
only when the model is committable.

The MigLayout body should be reviewed as a screen, not as a generated artifact.
The top rows contain identity fields. The address group spans both columns
because it has multiple related fields. The contacts table grows vertically.
The validation summary sits above the button row. The button row adapts
dialog commands and uses tags. The root panel is the component passed to
\api{FormDialog} for active-editor commit. Nothing about this layout requires
the controller to know field names.

The controller is created after the form model and root panel exist. Dirty
state may be a composition of field dirty flags, contact-table dirty state, and
pending uploaded documents. Cancel confirmation may call the application's
standard confirmation service. If the editor supports Apply, the surrounding
presenter must define what Apply does: update a preview, save a draft, or call
a repository. If Apply calls a repository, the successful Apply should move the
dirty baseline only after the repository confirms success, or the workflow
should make clear that Apply is a local preview rather than a save.

Validation focus requires a resolver. Name and credit fields can map through
component keys. Address problems may select an address section or tab. Contact
table problems translate model rows and columns to view coordinates because
the table may be sorted or filtered. The uniqueness problem for a customer code
targets the code field even though it was produced by a server lookup. The
resolver contains this navigation. The validation rules remain in the model or
validation service.

Background uniqueness validation needs stale-result suppression. Each code
change increments a generation. The submitted task captures the generation and
the candidate code. When the task returns, the presenter checks that the
dialog lifecycle is still open and the generation still matches. Only then
does it update the problem set, pending state, validation summary, and command
state. If the dialog has closed, the result is ignored. If the user changed the
code, the result is stale. This rule is worth writing in tests because it is
otherwise invisible until users see flickering or false errors.

Keyboard review checks that Enter from normal fields invokes OK, Enter inside
the contacts table respects table editing, Escape routes to dirty cancel, and
the close box follows the application's chosen cancel or cleanup policy. The
root pane binding handles the shared Escape/default-button part. Custom
components still need local input-map review. The goal is not to force every
component to obey the dialog; it is to ensure that when a modal decision is
made, it uses the same command path.

Lifecycle review checks that field bindings, table bindings, validation
binding, keyboard binding, behavior scope, task service, and any detachable
presenter state are registered. Closing the lifecycle should restore root-pane
state, component decorations, listeners, and task ownership before the
controller resets the form. A screenshot test can render the body in valid,
invalid, dirty, and busy states. A model test can assert the result record and
problem set. An EDT test can assert command behavior. A native-shell smoke test
can remain small.

\begin{center}
\textbf{Dialog review table.}
\end{center}

\noindent\begin{tabularx}{\linewidth}{@{}p{0.32\linewidth}X@{}}
\toprule
Observation & Review question\\
\midrule
OK loses the final table edit & Is active editor commit in the OK and Apply path?\\
OK closes with invalid input & Does \api{isCommittable} reflect parse, validation, and pending required checks?\\
Escape discards edits immediately & Is Escape routed to the same cancel command as the Cancel button?\\
Cancel asks during shutdown & Is user cancellation being confused with lifecycle cleanup?\\
Apply closes the dialog & Is Apply accidentally using OK semantics?\\
Apply then Cancel loses applied work & Is the dirty baseline moved after successful Apply, and does Cancel return the last applied result?\\
Revert behaves like Cancel & Is Revert backed by a real pre-dialog restoration policy rather than ordinary reset?\\
Validation summary changes but focus does not & Is a \api{ProblemFocusResolver} installed for structured targets?\\
Child problem opens the wrong tab & Are child focus resolvers composed with the same order and reveal policy as the composite session?\\
Problem appears on the wrong table row & Are model coordinates translated to view coordinates after sorting or filtering?\\
Old async validation result appears & Is generation or cancellation suppressing stale task results?\\
Dialog leaks listeners & Are bindings and task services registered with \api{FormDialogLifecycle}?\\
\bottomrule
\end{tabularx}

This review table is intentionally practical. Dialog defects are often
reported as anecdotes: ``sometimes OK ignores my edit'', ``Escape lost my
work'', ``the error points to the wrong row'', ``the dialog still updates after
I closed it.'' Each anecdote corresponds to a missing named rule. Naming the
rules lets a team review dialogs systematically instead of debugging each
report as a new mystery.

\section{Exercises and checklist}

\begin{exercisebox}
\begin{enumerate}
\item Find a dialog OK listener in an existing Swing program and list whether
it handles active editor commit, validation, focus recovery, result creation,
dirty cancel, keyboard consistency, and lifecycle cleanup.
\item Replace a direct OK button listener with a \api{FormDialog} command while
keeping the existing panel and layout.
\item Add a \api{ProblemFocusResolver} for a form with two fields and one
tabbed section. Write a test that proves the tab is selected before focus is
requested.
\item Write a dirty-cancel confirmation test that returns false and verifies
the controller remains open and the form is not reset.
\item Build a MigLayout dialog body with a validation summary, a growing table,
and a tagged button row. Render it with the screenshot harness at two sizes.
\item Add a rejecting table cell editor to a controller test and prove that OK
does not close while the editor rejects \api{stopCellEditing}.
\item Add Apply to a settings dialog and state the dirty-baseline policy in a
test.
\item Add asynchronous validation for a code field using a generation counter.
Prove that a stale result cannot publish a current problem.
\item Register a validation binding, keyboard binding, behavior scope, and task
service in \api{FormDialogLifecycle}. Prove that cleanup closes them before the
controller closes.
\item Review a generated form companion and decide which parts should be
generated facts and which parts should remain handwritten MigLayout.
\end{enumerate}
\end{exercisebox}

\begin{center}
\textbf{Dialog checklist.}
\end{center}

\begin{itemize}
\item The accepted result type is explicit and non-null.
\item Durable edit state lives in a form model, not only in widgets.
\item OK and Apply commit active editors before reading the form result.
\item OK is terminal only after successful commit and committability checks.
\item Apply is non-terminal and has a stated dirty-baseline policy.
\item Apply followed by Cancel preserves the last successfully applied result.
\item Revert is offered only when a real pre-dialog restoration policy exists.
\item Cancel uses dirty state and confirmation only for user cancellation.
\item Programmatic close bypasses dirty confirmation and performs cleanup.
\item Buttons, Enter, Escape, and tests invoke shared dialog commands.
\item Validation problems have stable codes, severities, and targets.
\item A focus resolver maps structured problem targets to Swing recovery.
\item Composite form sessions aggregate child problems before parent validation.
\item Generated child presentation models are composed, not committed.
\item MigLayout expresses the dialog body directly and remains reviewable.
\item Generated code supplies stable facts without hiding product layout.
\item Background work has cancellation, stale-result suppression, and EDT handoff.
\item Bindings, behaviors, task services, and detachables belong to a lifecycle.
\item The chapter examples are short class-name references, not long printed paths.
\end{itemize}

Dialog code should be boring in the best sense. The user may see a rich form,
careful validation, responsive background checks, and polished keyboard
behavior, but the implementation should have a small number of named moving
parts. A form model owns state. A panel owns geometry. Bindings connect state
and widgets. A controller owns the modal decision. A shell owns the native
window. A lifecycle owns temporary relationships. Once these parts are visible,
there is little room left for the traditional dialog tangle.
